\documentclass[12pt]{article}
\usepackage[T1]{fontenc}
\usepackage[utf8]{inputenc}
\usepackage{lmodern}
\usepackage{textcomp}
\usepackage{enumitem}
\usepackage{geometry}
\geometry{margin=1in}
\begin{document}
\begin{center}
Answer Key - Chemistry - Graduate Level Assessment 16 \
\small (Answers are ordered exactly as in the question paper.)
\end{center}

\section*{Multiple Choice}
\begin{enumerate}[label=\arabic*.]
\item \textbf{Answer:} A
\\ \textit{Options:} A) 0.185 $\mathrm{nm}$; B) 0.110 $\mathrm{nm}$; C) 0.125 $\mathrm{nm}$; D) 0.200 $\mathrm{nm}$; E) 0.095 $\mathrm{nm}$
\\ \textit{Note:} The correct answer of 0.185 nm is derived from the van der Waals radius for chlorine, which provides an estimate of the size of the $\mathrm{Cl}_2$ molecule when considered as a spherical entity, reflecting the effective distance at which the molecule's electron cloud interacts with other molecules.
\item \textbf{Answer:} A
\\ \textit{Options:} A) 6.75 \ensuremath{\times} 10^-10 m; B) 8.22 \ensuremath{\times} 10^-11 m; C) 2.15 \ensuremath{\times} 10^-12 m; D) 1.885 \ensuremath{\times} $10^7$ m s^-1; E) 1.22 \ensuremath{\times} 10^-9 m
\\ \textit{Note:} The correct answer of 6.75 × 10^-10 m arises from applying the Heisenberg uncertainty principle, which relates the uncertainty in position to the uncertainty in momentum, with the electron's acceleration through a potential difference resulting in a defined kinetic energy and corresponding momentum uncertainty.
\item \textbf{Answer:} A
\\ \textit{Options:} A) 22.6 \ensuremath{\times} 10^\{-3\} $\mathrm{~m}^3 \mathrm{~mol}^{-1}$; B) 0.0226 $\mathrm{~m}^3 \mathrm{~mol}^{-1}$; C) 2.26 \ensuremath{\times} 10^\{-6\} $\mathrm{~m}^3 \mathrm{~mol}^{-1}$; D) 226 \ensuremath{\times} 10^\{-6\} $\mathrm{~m}^3 \mathrm{~mol}^{-1}$; E) 2.26 \ensuremath{\times} 10^\{-2\} $\mathrm{~m}^3 \mathrm{~mol}^{-1}$
\\ \textit{Note:} The correct answer is 22.6 × 10^\{-3\} $m^3$ mol^\{-1\} because the van der Waals parameter $b$ represents volume per mole, and converting 0.0226 $dm^3$ to cubic meters requires multiplying by 10^\{-3\} to account for the factor that 1 $dm^3$ equals 10^\{-3\} $m^3$.
\item \textbf{Answer:} E
\\ \textit{Options:} A) Proteins determine the rate of photosynthesis in plants.; B) Proteins are used for insulation in cold environments; C) Proteins are the main constituent of cell walls in bacteria.; D) Proteins are responsible for coloration of organisms; E) Proteins serve as structural material and biological regulators.
\\ \textit{Note:} Proteins are essential to biological systems because they not only provide structural support to cells and tissues but also function as enzymes and hormones that regulate various biochemical processes.
\item \textbf{Answer:} A
\\ \textit{Options:} A) 0.311; B) 0.5; C) 1.00; D) 0.05; E) 0.10
\\ \textit{Note:} The mole fraction of sodium hydroxide is calculated by determining the number of moles of sodium hydroxide and water in a solution with 50\% weight by using their molar masses to convert the weight percent into moles, leading to a resultant mole fraction of approximately 0.311.
\end{enumerate}

\section*{True / False}
\begin{enumerate}[label=\arabic*.]
\setcounter{enumi}{5}
\item \textbf{Answer:} False
\\ \textit{Options:} A) True; B) False
\\ \textit{Note:} The statement is false because the resulting pH after adding NaOH to HCl depends on the relative concentrations of the acid and base, and if NaOH is in excess, the solution will have a pH greater than 7, indicating it is basic rather than neutral.
\item \textbf{Answer:} True
\\ \textit{Options:} A) True; B) False
\\ \textit{Note:} The statement is true because a 0.500 N solution of sulfuric acid (H$_{2}$SO$_{4}$) indicates a normality that accounts for the two hydrogen ions contributed by each molecule of H$_{2}$SO$_{4}$, resulting in a total mass of 102.5 grams in 3.00 liters of solution.
\item \textbf{Answer:} True
\\ \textit{Options:} A) True; B) False
\\ \textit{Note:} The statement is true because the infrared absorption frequency is inversely related to the square root of the reduced mass of the diatomic molecules, and since deuterium (D) has a greater mass than hydrogen (H), the reduced mass of DCl is greater than that of HCl, resulting in a lower absorption frequency for DCl.
\item \textbf{Answer:} True
\\ \textit{Options:} A) True; B) False
\\ \textit{Note:} The maximum kinetic energy of the β^- particle emitted during the decay of $He^6$ is calculated based on the mass difference between the parent and daughter nuclei, resulting in a value of 7.51 MeV due to the specific decay process and energy conservation principles involved.
\item \textbf{Answer:} True
\\ \textit{Options:} A) True; B) False
\\ \textit{Note:} The statement is true because the mole fraction, defined as the ratio of the number of moles of a component to the total number of moles in the solution, correctly indicates that in this case, H₂SO₄ constitutes 10\% of the total moles while H₂O constitutes 90\%.
\end{enumerate}

\section*{Long Form Response}
\begin{enumerate}[label=\arabic*.]
\setcounter{enumi}{10}
\item \textbf{Answer:} To calculate the work done by the gas during its expansion, we can use the formula \( W = -P_{\text{ext}} \Delta V \). The initial volume is 2.0 liters, and since the gas expands against an external pressure of 0.80 atm, we first need to find the change in volume, which is equal to the final volume minus the initial volume. Assuming the gas expands to 3.0 liters, the change in volume (\( \Delta V \)) is 1.0 liter. Thus, the work done is \( W = -0.80 \, \text{atm} \times 1.0 \, \text{L} = -0.80 \, \text{L atm} \). However, if the gas expands fully to 2.0 liters, the work done would be \( W = -2.4 \, \text{atm} \times 2.0 \, \text{L} = -4.8 \, \text{L atm} \), indicating that the work done by the system is -3.2 L atm when the expansion is considered against 0.80 atm. In thermodynamic terms, this negative sign signifies that work is done by the system on the surroundings, indicating an energy transfer that can affect both the internal energy and enthalpy of the gas.
\\ \textit{Note:} correct because it accurately applies the formula for work done during gas expansion, taking into account the external pressure and the change in volume, and correctly calculates the work as negative, indicating work done by the system on the surroundings.
\item \textbf{Answer:} The relationship between pH, pKa, and chemical shift in a solution of sodium phosphate is critical for understanding protonation states of phosphoric acid species. At the pKa of H$_{2}$PO$_{4}$-, which is approximately 7.2, the equilibrium between H$_{2}$PO$_{4}$- and HPO $4^2$- is established, influencing the protonation state and resulting chemical environment of phosphorus in the solution. At this pKa, the expected chemical shift for the phosphorus signal is around 6.10 ppm, indicating a distinct resonance that reflects the balance of protonation. This shift is a result of the chemical environment influenced by the acidic and basic forms of the phosphate species, illustrating how pH impacts the NMR chemical shifts of the molecules in solution. Understanding this interplay is essential for interpreting NMR spectra in biochemical and chemical contexts involving phosphate.
\\ \textit{Note:} correct because it accurately describes how the pH and pKa determine the protonation states of H$_{2}$PO$_{4}$- and HPO 4²-, which in turn influence the chemical environment and the resulting chemical shift of phosphorus in a sodium phosphate solution, particularly noting the expected shift at the pKa value of approximately 7.2.
\end{enumerate}

\vfill
\noindent\textit{End of Answer Key}
\end{document}